% Editable question content; compile CPE333_Midterm Exam.tex in this directory.
\sectiontitle{Part I: Multiple Choice (20 points)}
\begingroup\small
\noindent Circle the one correct answer in each question.\par\smallskip
\begin{questions}\setlength{\itemsep}{0.08cm}\setcounter{question}{0}
\question What is a trap in an operating system?
\begin{verticalchoices}
\item A trap is a system call that provides OS level service
\item A trap is a service that user can capture output from a process
\item A trap is a system call that captures users in OS
\item A trap is a way to capture bugs
\item There is no trap in OS
\end{verticalchoices}
\question What protection risk arises if user mode and kernel mode are not separated?
\begin{verticalchoices}
\item Process cannot be kicked of processor
\item OS cannot access disks or external devices
\item Process may be able to access other process data
\item OS can prevent process from directly access devices
\item Process cannot create child process
\end{verticalchoices}
\question Which set gives the four components of a process in the model used in class?
\begin{verticalchoices}
\item Program
\item Data state
\item Process status
\item Control state
\item All of the choices are process components
\end{verticalchoices}
\question Which OS goal is specifically intended to protect one application from another?
\begin{verticalchoices}
\item Performance
\item Isolation
\item Reliability
\item Mobility
\item All of the issues is for protection
\end{verticalchoices}
\question Which low-level hardware mechanism is needed in a preemptive system to switch a CPU-bound process?
\begin{verticalchoices}
\item MMU
\item Timer interrupt
\item Cache
\item TLB
\item None of them
\end{verticalchoices}
\question Which part of a process is stored in kernel space?
\begin{verticalchoices}
\item Text segment
\item Data segment
\item Stack segment
\item Process descriptor
\item All of them are not in kernel
\end{verticalchoices}
\end{questions}\vfill\newpage

\begin{questions}\setlength{\itemsep}{0.08cm}\setcounter{question}{6}
\question In a conventional process address space, which region grows in the opposite direction from the heap?
\begin{verticalchoices}
\item Text
\item Heap
\item Stack
\item Program
\item They are all growing in the same direction
\end{verticalchoices}
\question When would a process in the ready state be moved directly to the blocked state?
\begin{verticalchoices}
\item When it read a file
\item When it read from memory
\item When timer interrupt occurs
\item Cannot happen
\item Both when read from file and memory
\end{verticalchoices}
\question Which event can move a process into the ready state?
\begin{verticalchoices}
\item Timer interrupt occurs
\item I/O transfer is completed
\item A process is created
\item All of the events
\item None of the events
\end{verticalchoices}
\question What does it mean for a process to be a zombie?
\begin{verticalchoices}
\item Process is running on processor and cannot be kicked off
\item Process is waiting for I/O indefinitely
\item Process is ready but has never be select to run
\item Process is terminated but cannot be removed from the system
\item There is no such thing as zombie in OS
\end{verticalchoices}
\question Which process-control block stores a process's state and saved CPU context?
\begin{verticalchoices}
\item PCB
\item TLB
\item Stack
\item MMU
\item Queue
\end{verticalchoices}
\question Which situation can leave a zombie process in the system?
\begin{verticalchoices}
\item Parent exists before child returns value
\item Parent creates more than one child
\item Child also creates a child process
\item Parent waits for child
\item Child executes other program
\end{verticalchoices}
\end{questions}\vfill\newpage

\begin{questions}\setlength{\itemsep}{0.08cm}\setcounter{question}{12}
\question In the lab, when a parent waits for the child before printing, what will always happen?
\begin{verticalchoices}
\item[] (a) Child will always print after parent
\item[] (b) A zombie will always happen
\item[] (c) Child will always start from beginning of main function
\item[] (d) A segmentation fault will always occur
\item[] (e) Parent will always print after child
\end{verticalchoices}
\question The following C program:
\begin{lstlisting}[language=C,basicstyle={\normalfont},columns=fullflexible,keepspaces=true,showstringspaces=false,aboveskip=0pt,belowskip=0pt]
main()
{
    fork();
    fork();
    printf("hello");
}
\end{lstlisting}
How may time the word "hello" will be printed
\begin{verticalchoices}
\item Only once
\item Twice
\item Three times
\item Four times
\item Eight times
\end{verticalchoices}
\question Which events can cause a running process to give up the CPU?
\begin{verticalchoices}
\item Timer interrupt
\item Calling I/O transfer
\item Calling wait
\item Calling yield
\item All of the events
\end{verticalchoices}
\question Which scheduling algorithm executes jobs to completion in the order of arrival?
\begin{verticalchoices}
\item FIFO
\item SJF
\item RR
\item Both FIFO and RR
\item Both FIFO and SJF
\end{verticalchoices}
\question In the C storage-class model covered in class, which class requests register storage for a local variable?
\begin{verticalchoices}
\item Auto
\item Register
\item Extern
\item Static
\item They are executed at the same speed by CPU
\end{verticalchoices}
\end{questions}\vfill\newpage

\begin{questions}\setlength{\itemsep}{0.08cm}\setcounter{question}{17}
\question Which statement about partition allocation is true?
\begin{verticalchoices}
\item Single partition allocation is more efficient that multiple partition allocation in terms of memory usage
\item Flexible-size partition does not suffer external fragmentation
\item Fixed-size partition suffers from internal fragmentation
\item To solve internal fragmentation, OS needs to do memory coalescing and compaction
\item None of them is true
\end{verticalchoices}
\question What does the short-term process scheduler do?
\begin{verticalchoices}
\item Select program to load into memory
\item Select process to run
\item Select process to remove from memory
\item Select memory to assign to a process
\item Select process to put to memory
\end{verticalchoices}
\question What does a page-replacement algorithm do when a frame is needed?
\begin{verticalchoices}
\item Select a page to load into memory
\item Select a page to run
\item Select a page to remove from memory
\item Select memory to assign to a page
\item Select a page to replace a process
\end{verticalchoices}
\question What is not used in address translation?
\begin{verticalchoices}
\item MMU
\item Page table
\item TLB
\item Segment table
\item All of them are used in address translation
\end{verticalchoices}
\question Which option names the fragmentation problem that can occur with contiguous segment allocation?
\begin{verticalchoices}
\item Harder to translate from VA to PA
\item Cannot protect memory from other process
\item Suffer from external fragmentation
\item Use too much memory to keep map of the address translation
\item There is no disadvantage
\end{verticalchoices}
\question What does demand paging mean?
\begin{verticalchoices}
\item Page will be loaded when it is accesses only
\item Page is always in memory for on demand access
\item Page is pre-loaded on user demand
\item Page is pre-loaded by probability of the on demand usage
\item There is no such thing as on demand paging in OS
\end{verticalchoices}
\end{questions}\vfill\newpage

\begin{questions}\setlength{\itemsep}{0.08cm}\setcounter{question}{23}
\question Which step does NOT occur while handling a valid page fault?
\begin{verticalchoices}
\item Save process context
\item Free a frame if there is no free frame
\item Locate and read page from disk
\item Update page table
\item All events happen
\end{verticalchoices}
\question Which segment-table field stores a segment's starting physical address?
\begin{verticalchoices}
\item Base
\item Bound
\item Direction
\item Valid
\item All choices are a part of segment table
\end{verticalchoices}
\question With a fixed virtual address-space size, what is one effect of using a larger page size?
\begin{verticalchoices}
\item More problem with external fragmentation
\item Smaller TLB reaches
\item Smaller page table
\item More process can be created
\item There is no effect from the size of page
\end{verticalchoices}
\question Which single-queue policy rotates among all ready jobs, giving each a time slice in turn?
\begin{verticalchoices}
\item FIFO
\item STCF
\item RR
\item MLFQ
\item All of the schedulers have the same fairness
\end{verticalchoices}
\question Which scheduler does not need timer interrupt?
\begin{verticalchoices}
\item FIFO
\item STCF
\item RR
\item MLPQ
\item All of the schedulers have the same fairness
\end{verticalchoices}
\question Where does virtual memory save the image of memory?
\begin{verticalchoices}
\item On virtual storage
\item On a virtual cloud
\item On virtual private network
\item On kernel of the OS
\item On secondary storage
\end{verticalchoices}
\end{questions}\vfill\newpage

\begin{questions}\setlength{\itemsep}{0.08cm}\setcounter{question}{29}
\question What is the response time of a process?
\begin{verticalchoices}
\item[] (a) Time that process starts running until finish
\item[] (b) Time that process enters system until finish
\item[] (c) Time that process enters system until running
\item[] (d) Time that process runs until requests for I/O
\item[] (e) Time that process runs at each cycle
\end{verticalchoices}
\question What is the turnaround time of a process?
\begin{verticalchoices}
\item[(a)] Time that process starts running until finish
\item[(b)] Time that process enters system until finish
\item[(c)] Time that process enters system until running
\item[(d)] Time that process runs until requests for I/O
\item[(e)] Time that process runs at each cycle
\end{verticalchoices}
\question What operation may be needed to address external fragmentation from flexible-size partitions?
\begin{verticalchoices}
\item[(a)] Limit degree of multitasking
\item[(b)] Internal fragmentation
\item[(c)] Address translation
\item[(d)] Need compacting
\item[(e)] There is no problem
\end{verticalchoices}
\question With a single-level page table and fixed page size, what is a cost of a larger virtual address space?
\begin{verticalchoices}
\item[(a)] Need more memory to store page table
\item[(b)] Internal fragmentation
\item[(c)] External fragmentation
\item[(d)] TLB reach is smaller
\item[(e)] There is no problem causes by larger address space
\end{verticalchoices}
\question How can a multilevel page table save memory for a sparse address space?
\begin{verticalchoices}
\item[(a)] Table size is reduced
\item[(b)] Address space is reduced
\item[(c)] Page size is increased
\item[(d)] All of the second level pages are empty
\item[(e)] There is no need to store all the empty second level pages
\end{verticalchoices}
\question In the simplified state model, which state does a newly created process enter after admission?
\begin{verticalchoices}
\item[(a)] ready
\item[(b)] running
\item[(c)] blocked
\item[(d)] terminated
\item[(e)] empty
\end{verticalchoices}
\end{questions}\vfill\newpage

\begin{questions}\setlength{\itemsep}{0.08cm}\setcounter{question}{35}
\question Which state is a process in while it is executing on the CPU?
\begin{verticalchoices}
\item ready
\item running
\item blocked
\item terminated
\item empty
\end{verticalchoices}
\question What fairness risk can SJF create if shorter jobs keep arriving?
\begin{verticalchoices}
\item Suffer convoy effect
\item Highest average waiting time
\item Highest average response time
\item Highest average turn around time
\item Starvation
\end{verticalchoices}
\question Which scheduling metric is Round Robin primarily designed to improve?
\begin{verticalchoices}
\item Lowest average turn around time
\item Lowest average waiting time
\item starvation could occur
\item Lowest average response time
\item highest throughput
\end{verticalchoices}
\question Which statement about paging is correct?
\begin{verticalchoices}
\item Page and page frame can be different in size
\item The smaller the page size is, the smaller the page table is
\item Larger page transfer as fast as smaller page
\item Page is in logical memory, page frame is in physical memory
\item Demand paging is when a program can demand preloading of a page
\end{verticalchoices}
\question Which page-replacement policy is an ideal benchmark but cannot be implemented exactly without knowing future references?
\begin{verticalchoices}
\item Optimal
\item FIFO
\item LRU
\item LFU
\item All of the answers can be implemented
\end{verticalchoices}
\question When does a page fault occur?
\begin{verticalchoices}
\item When a process access a page that is not valid in page table
\item When a process access a page that is not in logical memory
\item When a process access a page that is updated in memory
\item When a process access a page not updated in memory
\item All of the answers cause page fault to occur
\end{verticalchoices}
\end{questions}\endgroup\vfill\newpage

\sectiontitle{Part II: Short Answer (10 points)}
\begingroup\small
\noindent Answer all questions briefly in the spaces provided. (1 point each.)\par\smallskip
\begin{questions}\setlength{\itemsep}{0.12cm}
\question Write the full names of MLFQ and STCF.\par\vspace{1.40cm}
\question What is the difference between a process and a program?\par\vspace{1.40cm}
\question Expand MMU and state its address-translation function.\par\vspace{1.40cm}
\question Give two examples of CPU state saved during a context switch.\par\vspace{1.40cm}
\question Give two types of information stored in a PCB, other than CPU register values.\par\vspace{1.40cm}
\question What is an exited child process called before its living parent reaps its exit status?\par\vspace{1.40cm}
\question Give one example each of an I/O-intensive task and a CPU-intensive task.\par\vspace{1.40cm}
\question What is thrashing in a virtual-memory system?\par\vspace{1.40cm}
\question With base = 16 KiB and bound = 4 KiB, is virtual byte address 4200 valid? Why?\par\vspace{1.40cm}
\question Expand TLB and explain why a TLB miss need not cause a page fault.\par\vspace{1.40cm}
\end{questions}\endgroup\vfill\newpage

\sectiontitle{Part III: Written Questions (20 points)}
\begingroup\small
\noindent Answer all three questions in the answer booklet. Show your work.\par\smallskip
\begin{questions}
\question \textbf{(8 points)} Consider the following set of processes (jobs), where the run time of the process is given in milliseconds (ms):
\begin{center}
\renewcommand{\arraystretch}{1.0}
\begin{tabular}{c|c|c}
Process & Run Time & Priority \\\hline
$P_1$ & 10 & 3 \\
$P_2$ & 1 & 1 \\
$P_3$ & 2 & 3 \\
$P_4$ & 1 & 4 \\
$P_5$ & 5 & 2
\end{tabular}
\end{center}
The processes are assumed to arrive in the order $P_1$, $P_2$, $P_3$, $P_4$, $P_5$, all at time 0.
\begin{parts}
\part Draw three diagrams illustrating the execution of these processes using FIFO, SJF, and RR scheduling with time slice = 1 ms.
\part What is the Turnaround time of each process for each of the schedulers in \mbox{part (a)?}
\part What is the Response time of each process for each of the schedulers in \mbox{part (a)?}
\part Which of the schedulers results in (i) the minimum average Turnaround time, (ii) the minimum average Response time (taking the average over all processes)?
\end{parts}

\question \textbf{(4 points)} Explain how a commercial bank handles customer during rush hours with the concept of the Multi-level Feedback Queue (MLFQ), in your opinion. Describe the algorithms and policies used and provide suggestions on how to improve the process with the goal of optimizing the average turnaround time of customers.

\question \textbf{(8 points)} Draw three tables show the 3 Frames in physical memory using FIFO, LRU and OPT for page replacement policy (separate table for both policies) and find the Effective Access Time (EAT) of each policies when references strings (list of needed pages) as below:
\begin{center}
7, 0, 1, 3, 0, 3, 2, 4, 2, 7, 0, 3, 2, 1, 2, 0, 1, 4, 0, 1 (20 needed pages)
\end{center}
Assume that the memory access time is 10 ns and the disk access time is 10 ms.
\end{questions}
\endgroup\vfill\newpage
% Pages 11 and 12 are intentionally blank scratch pages. Written answers still
% belong in the separate answer booklet, as specified on the unchanged cover.
\null\vfill\newpage
\null\vfill
